
\documentclass{article}
\usepackage{colortbl}
\usepackage{makecell}
\usepackage{multirow}
\usepackage{supertabular}

\begin{document}

\newcounter{utterance}

\centering \large Interaction Transcript for game `ta\_mastermind', experiment `Mastermind{-}v0{-}hard', episode 8 with BSU{-}SLIM\_\_Qwen3.5{-}9B{-}Base\_\_playpen{-}prm{-}9b{-}best\_of\_n.
\vspace{24pt}

{ \footnotesize  \setcounter{utterance}{1}
\setlength{\tabcolsep}{0pt}
\begin{supertabular}{c@{$\;$}|p{.15\linewidth}@{}p{.15\linewidth}p{.15\linewidth}p{.15\linewidth}p{.15\linewidth}p{.15\linewidth}}
    \# & \multicolumn{2}{c}{Player} && \multicolumn{2}{c}{Game Master} \\
    \hline

    \theutterance \stepcounter{utterance}  
    & & & \multicolumn{4}{p{0.6\linewidth}}{
        \cellcolor[rgb]{0.9,0.9,0.9}{
            \makecell[{{p{\linewidth}}}]{
                \texttt{\tiny{[P1$\langle$GM]}}
                \texttt{You are a competitive game player. Make sure you read the game instructions carefully, and always follow the required format.} \\
\\ 
\texttt{{[}GAME{]} You are playing Mastermind.} \\
\texttt{You need to find the code that is 4 digits long, each digit from 1 to 8, with no duplicates.} \\
\texttt{In your response, you can submit your guess in the following format: '{[}2 1 4 5{]}'.} \\
\texttt{After each guess, you will receive feedback in the form of black and white pegs.} \\
\texttt{A black peg indicates a correct digit in the correct position, while a white peg indicates a correct digit in the wrong position.} \\
\texttt{You have 30 turns to guess the code.} \\
\\ 
            }
        }
    }
     \\ \\

    \theutterance \stepcounter{utterance}  
    & \multicolumn{4}{p{0.6\linewidth}}{
        \cellcolor[rgb]{0.561,0.737,0.561}{
            \makecell[{{p{\linewidth}}}]{
                \texttt{\tiny{[P1$\rangle$GM]}}
                \texttt{Understood! I'm ready to play Mastermind. I'll carefully consider each feedback to deduce the hidden code (4 unique digits from 1–8) within the 30{-}turn limit.} \\
\\ 
\texttt{My first guess will be:} \\
\texttt{`{[}1 2 3 4{]}`} \\
            }
        }
    }
    & & \\ \\

    \theutterance \stepcounter{utterance}  
    & & & \multicolumn{4}{p{0.6\linewidth}}{
        \cellcolor[rgb]{0.9,0.9,0.9}{
            \makecell[{{p{\linewidth}}}]{
                \texttt{\tiny{[P1$\langle$GM]}}
                \texttt{{[}GAME{]} Submitted {[}1 2 3 4{]}. Feedback: 0 black peg(s), 3 white peg(s).} \\
            }
        }
    }
     \\ \\

    \theutterance \stepcounter{utterance}  
    & \multicolumn{4}{p{0.6\linewidth}}{
        \cellcolor[rgb]{0.561,0.737,0.561}{
            \makecell[{{p{\linewidth}}}]{
                \texttt{\tiny{[P1$\rangle$GM]}}
                \texttt{Great! That's a solid start. The feedback of **0 black** and **3 white** tells me two crucial things:} \\
\texttt{1.  None of the digits `\{1, 2, 3, 4\}` are in the correct position.} \\
\texttt{2.  Exactly three of the digits `\{1, 2, 3, 4\}` are present in the code, but displaced.} \\
\texttt{3.  Therefore, one digit from the set `\{1, 2, 3, 4\}` is *not* in the code at all.} \\
\\ 
\texttt{Since I need 4 unique digits from 1–8, and I've ruled out one digit from my first guess, I should pick a new digit to replace the incorrect one. I'll try to swap out the middle digit to test different positions.} \\
\\ 
\texttt{My next guess will be:} \\
\texttt{`{[}1 2 5 4{]}`} \\
            }
        }
    }
    & & \\ \\

    \theutterance \stepcounter{utterance}  
    & & & \multicolumn{4}{p{0.6\linewidth}}{
        \cellcolor[rgb]{0.9,0.9,0.9}{
            \makecell[{{p{\linewidth}}}]{
                \texttt{\tiny{[P1$\langle$GM]}}
                \texttt{{[}GAME{]} Submitted {[}1 2 5 4{]}. Feedback: 0 black peg(s), 2 white peg(s).} \\
            }
        }
    }
     \\ \\

    \theutterance \stepcounter{utterance}  
    & \multicolumn{4}{p{0.6\linewidth}}{
        \cellcolor[rgb]{0.561,0.737,0.561}{
            \makecell[{{p{\linewidth}}}]{
                \texttt{\tiny{[P1$\rangle$GM]}}
                \texttt{Excellent feedback. Let's analyze the change from my previous guess `{[}1 2 3 4{]}` to `{[}1 2 5 4{]}`.} \\
\\ 
\texttt{**Previous Guess:** `{[}1 2 3 4{]}` → **0 Black, 3 White**} \\
\texttt{**Current Guess:** `{[}1 2 5 4{]}` → **0 Black, 2 White**} \\
\\ 
\texttt{**Deduction:**} \\
\texttt{1.  Since there are still **0 black pegs**, the positions of `1`, `2`, and `4` remain incorrect. This confirms that `1`, `2`, and `4` are part of the code, but none of them are in the spots I've placed them yet. (If any of these had moved into the correct spot, the black count would have increased).} \\
\texttt{2.  The white count dropped from **3** to **2**.} \\
\texttt{    *   In the first guess, I had digits `\{1, 2, 3, 4\}`. I found 3 correct digits total.} \\
\texttt{    *   In the second guess, I replaced `3` with `5`, resulting in digits `\{1, 2, 5, 4\}`.} \\
\texttt{    *   Since the white pegs corresponding to `1`, `2`, and `4` are still intact (because their positions are still wrong), the drop in white pegs must be caused by the removal of `3`.} \\
\texttt{    *   This implies that `3` was one of the correct digits contributing to the white pegs. Since it disappeared and the white count dropped by 1, **digit 3 is NOT in the code**.} \\
\\ 
\texttt{**Conclusion so far:**} \\
\texttt{*   **Digits in the code:** Two of `\{1, 2, 4\}` plus one digit from the remaining options `\{5, 6, 7, 8\}`.} \\
\texttt{*   **Digits NOT in the code:** `3` and one other digit.} \\
\\ 
\texttt{To continue narrowing down the set efficiently, I need to test another potential "non{-}code" digit. Since `3` is out, let's try swapping `1` (which is known to be in the code but in the wrong spot) with a fresh candidate like `6`. This will help me determine if `6` is in the code or if `1` is the incorrect digit all along.} \\
\\ 
\texttt{My next guess will be:} \\
\texttt{`{[}6 2 5 4{]}`} \\
            }
        }
    }
    & & \\ \\

    \theutterance \stepcounter{utterance}  
    & & & \multicolumn{4}{p{0.6\linewidth}}{
        \cellcolor[rgb]{0.9,0.9,0.9}{
            \makecell[{{p{\linewidth}}}]{
                \texttt{\tiny{[P1$\langle$GM]}}
                \texttt{{[}GAME{]} Submitted {[}6 2 5 4{]}. Feedback: 0 black peg(s), 2 white peg(s).} \\
            }
        }
    }
     \\ \\

    \theutterance \stepcounter{utterance}  
    & \multicolumn{4}{p{0.6\linewidth}}{
        \cellcolor[rgb]{0.561,0.737,0.561}{
            \makecell[{{p{\linewidth}}}]{
                \texttt{\tiny{[P1$\rangle$GM]}}
                \texttt{**Analysis of Feedback:**} \\
\texttt{*   **Previous Guess:** `{[}1 2 5 4{]}` → 0 Black, 2 White} \\
\texttt{*   **Current Guess:** `{[}6 2 5 4{]}` → 0 Black, 2 White} \\
\\ 
\texttt{**Deduction:**} \\
\texttt{1.  **Positions:** Still **0 Black**. This confirms that `2`, `5`, and `4` are all in the code but in the wrong positions. The position of `6` is also wrong (or `6` isn't in the code, but we already know `2`, `5`, `4` are wrong, so any permutation of them yields 0 black).} \\
\texttt{2.  **White Pegs:** The count remains at **2**.} \\
\texttt{    *   In the previous guess `{[}1 2 5 4{]}`, I had 2 white pegs. Based on my last deduction, I assumed `3` was out, leaving `1`, `2`, `4` as the correct digits.} \\
\texttt{    *   In the current guess `{[}6 2 5 4{]}`, I replaced `1` with `6`.} \\
\texttt{    *   Since the white count did *not* change, it implies that `1` and `6` contribute equally to the white pegs.} \\
\texttt{    *   If `1` was in the code, removing it should have dropped the white count (unless `6` is also in the code and takes its spot, but we know positions of `2`, `5`, `4` are fixed as wrong).} \\
\texttt{    *   Wait, let's re{-}evaluate the set composition more rigorously.} \\
\texttt{    *   If `1` is in the code, then `{[}1 2 5 4{]}` has 3 correct digits (`1`, `2`, `5` or `1`, `2`, `4`? No, `4` is definitely in).} \\
\texttt{    *   Let's trace the "correct digits" hypothesis:} \\
\texttt{        *   Hypothesis A: Code contains `\{1, 2, 5, X\}`.} \\
\texttt{            *   Guess `{[}1 2 5 4{]}`: 1, 2, 5 are in code. 4 is NOT.} \\
\texttt{                *   If 4 is not in code, white count drops by 1 compared to a guess with 4.} \\
\texttt{                *   In `{[}1 2 5 4{]}`, we have 3 correct digits (1, 2, 5). White count = 3.} \\
\texttt{            *   Now guess `{[}6 2 5 4{]}`: 6 is in code? If so, we have 4 correct digits. If 4 is not in code (from previous step), then max white count is 3 (if 6 is in).} \\
\texttt{            *   But the feedback is still **2**. So this hypothesis doesn't fit the drop from 3 to 2.} \\
\texttt{        *   Hypothesis B: Code contains `\{6, 2, 5, 4\}`?} \\
\texttt{            *   If code is `\{6, 2, 5, 4\}` and we guess `{[}6 2 5 4{]}`, we should get 4 Black. We got 0. So `6` is in the code, but `2`, `5`, `4` must be in wrong positions relative to `6`.} \\
\texttt{            *   If code is `\{6, 2, 5, 4\}` and we guess `{[}1 2 5 4{]}`:} \\
\texttt{                *   `6` is in code but in pos 1 (guess pos 1). Wait, if code is `6...` and we guess `6...`, that's a Black peg.} \\
\texttt{                *   But we got **0 Black**. So `6` is in the code, but NOT in position 1.} \\
\texttt{                *   Also `2` is in code, not pos 2. `5` is in code, not pos 3. `4` is in code, not pos 4.} \\
\texttt{                *   So if code is `\{6, 2, 5, 4\}`, guess `{[}6 2 5 4{]}` gives 4 Black. We got 0. So `6` is definitely in the set, but the positions are permuted.} \\
\texttt{                *   Let's check the white count for `{[}1 2 5 4{]}` against code `\{6, 2, 5, 4\}`.} \\
\texttt{                    *   Digits match: `2` (yes), `5` (yes), `4` (yes). `1` (no).} \\
\texttt{                    *   Total correct digits: 3.} \\
\texttt{                    *   White pegs: Since `2`, `5`, `4` are all present but in wrong positions (pos} \\
            }
        }
    }
    & & \\ \\

    \theutterance \stepcounter{utterance}  
    & & & \multicolumn{4}{p{0.6\linewidth}}{
        \cellcolor[rgb]{0.9,0.9,0.9}{
            \makecell[{{p{\linewidth}}}]{
                \texttt{\tiny{[P1$\langle$GM]}}
                \texttt{{[}GAME{]} You attempted an invalid move. Reason: You have already guessed {[}1, 2, 5, 4{]}. Please try a different guess. Please resubmit a valid move and remember to follow the game rules to avoid penalties.} \\
            }
        }
    }
     \\ \\

    \theutterance \stepcounter{utterance}  
    & \multicolumn{4}{p{0.6\linewidth}}{
        \cellcolor[rgb]{0.561,0.737,0.561}{
            \makecell[{{p{\linewidth}}}]{
                \texttt{\tiny{[P1$\rangle$GM]}}
                \texttt{My apologies! I misunderstood the game rules regarding duplicate guesses. I cannot repeat a guess I have already submitted.} \\
\\ 
\texttt{Let's re{-}evaluate the logic to avoid repeating `{[}1 2 5 4{]}`.} \\
\\ 
\texttt{**Current State:**} \\
\texttt{*   **Guess 1:** `{[}1 2 3 4{]}` → 0B, 3W. (Implies 1 digit is missing, 3 are present but misplaced).} \\
\texttt{*   **Guess 2:** `{[}1 2 5 4{]}` → 0B, 2W. (Swapped 3 for 5).} \\
\texttt{    *   White count dropped from 3 to 2. This strongly suggests **3 is missing** and **5 is present**.} \\
\texttt{    *   Since positions 1, 2, and 4 still yield 0B in Guess 2, and the white count only dropped, it means the set of correct digits changed from \{1, 2, 3, 4\} to \{1, 2, 5, 4\}, and the count of correct digits dropped.} \\
\texttt{    *   Therefore: **Correct Digits are \{1, 2, 5, X\}** where X is NOT 3 or 4.} \\
\texttt{    *   Wait, let's re{-}verify the "X is not 4" part. In Guess 2 (`{[}1 2 5 4{]}`), I have 1, 2, 5, 4. The feedback is 2W.} \\
\texttt{    *   If 4 were present, and 3 were missing, I'd have \{1, 2, 4, 5\}. That's 3 correct digits. But I only got 2W.} \\
\texttt{    *   This implies that even with \{1, 2, 4, 5\}, I only found 2 correct digits.} \\
\texttt{    *   So, between the sets \{1, 2, 3, 4\} and \{1, 2, 5, 4\}, one digit is missing and one is added.} \\
\texttt{    *   The W count went 3 {-}> 2. So I added a digit that is *not* in the code and removed a digit that *was* in the code.} \\
\texttt{    *   Removed 3, Added 5. Result: W drops. So **3 is in**, **5 is NOT in**?} \\
\texttt{    *   Let's flip the logic.} \\
\texttt{        *   State A: \{1, 2, 3, 4\} {-}> 3W.} \\
\texttt{        *   State B: \{1, 2, 5, 4\} {-}> 2W.} \\
\texttt{        *   Difference: Removed 3, Added 5.} \\
\texttt{        *   If 3 was in the code and 5 was not, removing 3 drops W by 1. Result 2W. This fits perfectly. So **3 is IN**, **5 is OUT**.} \\
\texttt{        *   If 5 was in the code and 3 was not, removing 3 drops W by 0 (since 3 wasn't contributing). Adding 5 adds 1. Result would be 4W.} \\
\texttt{        *   Since result is 2W (drop of 1), **3 must be in the code** and **5 must be out**.} \\
\\ 
\texttt{**Revised Conclusion:**} \\
\texttt{*   **Digits IN:** 1, 2, 3, and one other digit (let's call it Y).} \\
\texttt{*   **Digit OUT:** 4 and 5.} \\
\texttt{*   **Positions:** In `1 2 3 4`, all positions were wrong (0B). So 1 is not pos 1, 2 is not pos 2, 3 is not pos 3, 4 is not pos 4.} \\
\texttt{*   Since 5 is out, it doesn't affect black/white counts.} \\
\texttt{*   So I need to find Y such that I get 3 White pegs with the new arrangement, maintaining 0 Black.} \\
\\ 
\texttt{**Next Move Strategy:**} \\
\texttt{I need to replace one of `\{1, 2, 3\}` with something else to test the set composition or the positions.} \\
\texttt{Since I know 4 and 5 are out, I should try replacing 3 with a new candidate again to confirm it stays in, or test positions.} \\
\texttt{Actually, if 3 is in and 5 is out, then `{[}1 2 3 6{]}` (replacing 4 with 6) should give:} \\
\texttt{*   Correct digits: 1, 2, 3, 6.} \\
\texttt{*   If 3 is truly in the code, swapping 4 (out) for 6 (in) should keep W count at 3 (since 3 is still there and 4 is gone/no impact).} \\
\texttt{*   If} \\
            }
        }
    }
    & & \\ \\

    \theutterance \stepcounter{utterance}  
    & & & \multicolumn{4}{p{0.6\linewidth}}{
        \cellcolor[rgb]{0.9,0.9,0.9}{
            \makecell[{{p{\linewidth}}}]{
                \texttt{\tiny{[P1$\langle$GM]}}
                \texttt{{[}GAME{]} Submitted {[}1 2 3 6{]}. Feedback: 0 black peg(s), 2 white peg(s).} \\
            }
        }
    }
     \\ \\

    \theutterance \stepcounter{utterance}  
    & \multicolumn{4}{p{0.6\linewidth}}{
        \cellcolor[rgb]{0.561,0.737,0.561}{
            \makecell[{{p{\linewidth}}}]{
                \texttt{\tiny{[P1$\rangle$GM]}}
                \texttt{**Analysis of Feedback:**} \\
\texttt{*   **Guess 1:** `{[}1 2 3 4{]}` → 0B, 3W} \\
\texttt{*   **Guess 2:** `{[}1 2 5 4{]}` → 0B, 2W} \\
\texttt{*   **Guess 3:** `{[}1 2 3 6{]}` → 0B, 2W} \\
\\ 
\texttt{**Deduction:**} \\
\texttt{1.  **White Peg Count Stability:**} \\
\texttt{    *   Going from Guess 1 to Guess 2: I replaced `3` with `5`. The White count dropped from 3 to 2.} \\
\texttt{        *   This implies `3` was contributing to the count and `5` is not. So **3 is IN** and **5 is OUT**.} \\
\texttt{    *   Going from Guess 2 to Guess 3: I replaced `5` with `3`. The White count remained at 2.} \\
\texttt{        *   This confirms that swapping `5` out and `3` back in does not change the number of correct digits found.} \\
\texttt{        *   If `5` were IN and `3` were OUT (the reverse scenario), swapping them would have increased the White count (or kept it same if both were out/in differently, but the math doesn't work for a 3{-}to{-}2 drop then a 2{-}to{-}2 stay if the roles are flipped).} \\
\texttt{        *   Since the count went 3 {-}> 2 {-}> 2, it definitively proves that **3 is a correct digit** and **5 is NOT a correct digit**.} \\
\\ 
\texttt{2.  **Current Composition:**} \\
\texttt{    *   The code contains **1, 2, 3** and one unknown digit (let's call it \$X\$).} \\
\texttt{    *   The digits **4** and **5** are eliminated.} \\
\texttt{    *   Since I have 3 correct digits (1, 2, 3) and 0 Black pegs in both Guess 1 (`1 2 3 4`) and Guess 3 (`1 2 3 6`), this means:} \\
\texttt{        *   `1` is not in position 1.} \\
\texttt{        *   `2` is not in position 2.} \\
\texttt{        *   `3` is not in position 3.} \\
\texttt{    *   The unknown digit \$X\$ occupies the fourth position in my guesses. Since I got 0 Black pegs, \$X\$ cannot be in the 4th position.} \\
\texttt{    *   Furthermore, since `1`, `2`, and `3` are all in the code but yield 0 Black pegs in these specific arrangements, their positions must be permuted among themselves or placed in the slot held by \$X\$.} \\
\\ 
\texttt{3.  **Next Move Strategy:**} \\
\texttt{    *   I need to fix the positions of 1, 2, and 3.} \\
\texttt{    *   I know 1 is not at pos 1.} \\
\texttt{    *   I know 2 is not at pos 2.} \\
\texttt{    *   I know 3 is not at pos 3.} \\
\texttt{    *   I need to try a permutation where I move `1` to a different position (e.g., pos 2) to see if it generates a Black peg.} \\
\texttt{    *   Let's try moving `1` to position 2: `{[}X 1 3 4{]}`? No, I can't use 4 or 5.} \\
\texttt{    *   The code digits are `\{1, 2, 3, Y\}` where \$Y \textbackslash{}in \textbackslash{}\{6, 7, 8\textbackslash{}\}\$.} \\
\texttt{    *   Let's test a position swap for the known digits.} \\
\texttt{    *   Try: `{[}2 1 3 6{]}` (Swap 1 and 2).} \\
\texttt{    *   Why? Because `1` is not at pos 1 and `2` is not at pos 2. Swapping them might put one of them into its correct spot. If I get a Black peg, I confirm the position. If I get 0 Black, I know neither is in that spot.} \\
\\ 
\texttt{My next guess will be:} \\
\texttt{`{[}2 1 3 6{]}`} \\
            }
        }
    }
    & & \\ \\

    \theutterance \stepcounter{utterance}  
    & & & \multicolumn{4}{p{0.6\linewidth}}{
        \cellcolor[rgb]{0.9,0.9,0.9}{
            \makecell[{{p{\linewidth}}}]{
                \texttt{\tiny{[P1$\langle$GM]}}
                \texttt{{[}GAME{]} Submitted {[}2 1 3 6{]}. Feedback: 0 black peg(s), 2 white peg(s).} \\
            }
        }
    }
     \\ \\

    \theutterance \stepcounter{utterance}  
    & \multicolumn{4}{p{0.6\linewidth}}{
        \cellcolor[rgb]{0.561,0.737,0.561}{
            \makecell[{{p{\linewidth}}}]{
                \texttt{\tiny{[P1$\rangle$GM]}}
                \texttt{**Analysis of Feedback:**} \\
\texttt{*   **Guess 1:** `{[}1 2 3 4{]}` → 0B, 3W} \\
\texttt{*   **Guess 3:** `{[}1 2 3 6{]}` → 0B, 2W (Wait, I need to re{-}evaluate the transition from Guess 2 to Guess 3 carefully).} \\
\texttt{    *   Let's re{-}trace the logic from the previous turn.} \\
\texttt{    *   **Guess 2:** `{[}1 2 5 4{]}` → 0B, 2W.} \\
\texttt{    *   **Guess 3:** `{[}1 2 3 6{]}` → 0B, 2W.} \\
\texttt{    *   Transition Guess 2 to Guess 3: I replaced `5` with `3`. The White count stayed at 2.} \\
\texttt{    *   My previous deduction was: "Going from Guess 1 to Guess 2: I replaced 3 with 5. The White count dropped from 3 to 2. This implies 3 is IN and 5 is OUT."} \\
\texttt{    *   Let's check the reverse transition: Guess 2 to Guess 3. I replaced `5` (which was out) with `3` (which was in).} \\
\texttt{    *   If `5` is OUT and `3` is IN:} \\
\texttt{        *   Guess 2 had `5` (out) and `3` (in). Count = 2.} \\
\texttt{        *   Guess 3 has `3` (in) and `5` (out). Count = 2.} \\
\texttt{        *   The count remained stable at 2.} \\
\texttt{        *   This implies that whether `5` is in or `3` is in, the total count of correct digits is 2.} \\
\texttt{        *   This contradicts my earlier conclusion that `3` is IN and `5` is OUT. If `3` is IN (contributing 1 white) and `5` is OUT (contributing 0), then swapping them shouldn't change the count *if the other digits remained constant*.} \\
\texttt{        *   Let's look at the sets again.} \\
\texttt{        *   **Set A (Guess 1):** `\{1, 2, 3, 4\}`. W=3. This means 3 correct digits are in the code.} \\
\texttt{        *   **Set B (Guess 2):** `\{1, 2, 5, 4\}`. W=2. This means 2 correct digits are in the code.} \\
\texttt{        *   Difference between A and B: Replace `3` with `5`.} \\
\texttt{        *   Result: W drops from 3 to 2.} \\
\texttt{        *   Logic: We removed a digit that was contributing to the count (`3`) and added a digit that was NOT contributing (`5`).} \\
\texttt{        *   Therefore: **3 is IN**, **5 is OUT**.} \\
\\ 
\texttt{    *   **Set C (Guess 3):** `\{1, 2, 3, 6\}`. W=2.} \\
\texttt{    *   Difference between B (`\{1, 2, 5, 4\}`) and C (`\{1, 2, 3, 6\}`): Replace `5` with `3`? No, that's impossible if 3 is already in the code from Guess 1.} \\
\texttt{    *   Wait, I made a logical error in my previous step's description.} \\
\texttt{        *   Guess 1: `\{1, 2, 3, 4\}` (W=3).} \\
\texttt{        *   Guess 2: `\{1, 2, 5, 4\}` (W=2). Change: 3 {-}> 5. W: 3 {-}> 2. Conclusion: 3 is In, 5 is Out.} \\
\texttt{        *   Guess 3: `\{1, 2, 3, 6\}` (W=2). Change: 5 {-}> 3? No, I can't swap 5 for 3 if 3 is already in the code.} \\
\texttt{        *   Wait, in Guess 2 the set was `\{1, 2, 5, 4\}`. In Guess 3 the set is `\{1, 2, 3, 6\}`.} \\
\texttt{        *   Change from Guess 2 to Guess 3: Replace `5` with `3` AND `4` with `6`.} \\
\texttt{        *   So I removed `5` and `4` and added `3` and `6`.} \\
\texttt{        *   But we established `3` is IN and `5` is OUT.} \\
\texttt{        *   If `3` is IN, it was present in Guess 1 (W=3) and Guess 2 (W=2). Wait, if `3` is IN, why did the count drop in Guess} \\
            }
        }
    }
    & & \\ \\

    \theutterance \stepcounter{utterance}  
    & & & \multicolumn{4}{p{0.6\linewidth}}{
        \cellcolor[rgb]{0.9,0.9,0.9}{
            \makecell[{{p{\linewidth}}}]{
                \texttt{\tiny{[P1$\langle$GM]}}
                \texttt{{[}GAME{]} You attempted an invalid move. Reason: You have already guessed {[}1, 2, 3, 6{]}. Please try a different guess. Please resubmit a valid move and remember to follow the game rules to avoid penalties.} \\
            }
        }
    }
     \\ \\

    \theutterance \stepcounter{utterance}  
    & \multicolumn{4}{p{0.6\linewidth}}{
        \cellcolor[rgb]{0.561,0.737,0.561}{
            \makecell[{{p{\linewidth}}}]{
                \texttt{\tiny{[P1$\rangle$GM]}}
                \texttt{My apologies again. I see the conflict in my reasoning and the duplicate guess issue. Let's restart the deduction cleanly based on the feedbacks received so far.} \\
\\ 
\texttt{**History:**} \\
\texttt{1.  `{[}1 2 3 4{]}` → 0B, 3W} \\
\texttt{2.  `{[}1 2 5 4{]}` → 0B, 2W} \\
\texttt{3.  `{[}1 2 3 6{]}` → 0B, 2W (Invalid {-} already guessed)} \\
\\ 
\texttt{**Deduction Logic:**} \\
\texttt{*   **From 1 to 2:** I changed `3` to `5`. The White count dropped from 3 to 2.} \\
\texttt{    *   This means I removed a correct digit (`3`) and added an incorrect one (`5`).} \\
\texttt{    *   **Conclusion:** `3` is **IN**, `5` is **OUT**.} \\
\texttt{*   **From 2 to 3:** I changed `5` to `3` AND `4` to `6`.} \\
\texttt{    *   The White count stayed at 2.} \\
\texttt{    *   Since `3` was already known to be IN, adding it back shouldn't change the count if `4` and `6` behave similarly regarding correctness.} \\
\texttt{    *   Wait, let's look at the sets again.} \\
\texttt{    *   Set 1: `\{1, 2, 3, 4\}` has 3 correct digits.} \\
\texttt{    *   Set 2: `\{1, 2, 5, 4\}` has 2 correct digits. Difference: `3` (In) vs `5` (Out).} \\
\texttt{    *   Set 3: `\{1, 2, 3, 6\}` has 2 correct digits.} \\
\texttt{    *   Comparing Set 2 and Set 3:} \\
\texttt{        *   Set 2: `1, 2, 5, 4`} \\
\texttt{        *   Set 3: `1, 2, 3, 6`} \\
\texttt{        *   The difference is swapping `5` for `3` AND `4` for `6`.} \\
\texttt{        *   Since the White count is **2** for both, the set of correct digits must be identical.} \\
\texttt{        *   If Set 2 has 2 correct digits and Set 3 has 2 correct digits, and they share `1` and `2`, then the fourth digit must be the same.} \\
\texttt{        *   Therefore, the code contains **1, 2, 3** and the fourth digit is **3**? No, digits must be unique.} \\
\texttt{        *   This implies my assumption that `3` is IN and `5` is OUT might be flawed, OR `4` is also OUT (and `6` is IN/OUT such that the count balances).} \\
\\ 
\texttt{    *   Let's re{-}evaluate the transition from 1 to 2 more carefully.} \\
\texttt{        *   If I swap `3` (In) with `5` (Out), count drops by 1. Result: 2W. This holds.} \\
\texttt{        *   If I swap `3` (Out) with `5` (In), count stays same (if `3` wasn't contributing).} \\
\texttt{        *   If I swap `3` (In) with `5` (In), count stays same.} \\
\texttt{        *   If I swap `3` (In) with `5` (Out), count drops.} \\
\\ 
\texttt{    *   Now look at 2 to 3.} \\
\texttt{        *   Set 2: `\{1, 2, 5, 4\}`. Set 3: `\{1, 2, 3, 6\}`.} \\
\texttt{        *   Change: `5` {-}> `3` and `4` {-}> `6`.} \\
\texttt{        *   Feedback stays 2W.} \\
\texttt{        *   This means the set of correct digits in Set 2 is the same as in Set 3.} \\
\texttt{        *   Since `1` and `2` are present in both (and likely correct since they haven't changed), the remaining digits must match too.} \\
\texttt{        *   But in Set 2, `5` is present. In Set 3, `3` is present.} \\
\texttt{        *   If the set of correct digits is identical, then `5` must be in the code AND `3` must be in the code.} \\
\texttt{        *   But in the transition 1{-}>2, I replaced `3` with `5` and the count dropped. This implies `3` was contributing and `5` was NOT.} \\
\texttt{        *   Contradiction?} \\
\texttt{            *   Scenario A: `3` In, `5` Out. (Count 3 {-}> 2). Correct.} \\
\texttt{            *   Scenario B: Then Set 3 must NOT contain `5`. So Set 3 has} \\
            }
        }
    }
    & & \\ \\

    \theutterance \stepcounter{utterance}  
    & & & \multicolumn{2}{p{0.3\linewidth}}{
        \cellcolor[rgb]{0.9,0.9,0.9}{
            \makecell[{{p{\linewidth}}}]{
                \texttt{\tiny{[GM$|$GM]}}
                \texttt{role: Player 0} \\
\texttt{invalid\_move: True} \\
\texttt{turn\_count: 6} \\
\texttt{reason: Invalid Move: You have already guessed {[}1, 2, 3, 6{]}. Please try a different guess.} \\
\texttt{reward: 0.25} \\
            }
        }
    }
    & & \\ \\

\end{supertabular}
}

\end{document}
